\SourceLeaf{532}$-11713$ und die durch Formen der resp. Fundamentalclassen dargestellten Zahlen wie $31\pdot\ 37\pdot\ 2$ beifüge, indem, wie in meiner Abhandlung Band 14 der Abhandlungen der Königlichen Gesellschaft der Wissenschaften zu Göttingen, als Fundamentalclassen solche Classen genommen werden, die in Vereinigung mit den Classen ihrer Perioden durch Composition jede eigentlich primitive Classe des Determinanten einmal und nur einmal hervorbringen.
\par\needspace{5\baselineskip}Es sind schon die Angaben fortgelassen:
\begin{longtable}{rcrl}
Centas 9 & G. IV $\mathinner{.\,.\,.}$ & $3\pdot,$ & $-827[21::3]$ \\
26 & IV & $14$ & $-2587[24::11]$ \\
26 & VIII & $6$ & $-2564[56::3]$ \\
91 & I & $111$ & $-9059[117::5]$ \\
120 & IV & $32$ & $-11956{}^{*}2^{*}[36\pdot2::11\pdot49]$ \\
1 & I & $2$ & $+37[3::3]$ \\
2 & I & $2$ & $+101[3::4]$ \\
\end{longtable}
\par\needspace{5\baselineskip}und hinzugefügt:
\begin{longtable}{rcrl}
Centas 9 & G. IV $\mathinner{.\,.\,.}$ & $3\dotstwo$ & $-828[6\pdot2::31\pdot23]$ \\
26 & IV & $14$ & $-2586[28\pdot2::7\pdot2]$ \\
26 & VIII & $6$ & $-2565[12\pdot2\pdot2::7\pdot2\pdot5]$ \\
91 & I & $117$ & $-9059[117::5]$ \\
120 & IV & $32$ & $-11966{}^{*}2^{*}[32\pdot4::5\pdot83]$ \\
1 & I & $3$ & $+37[3::3]$ \\
2 & I & $3$ & $+101[3::4]$ \\
\end{longtable}
\par Bei der Tafel für Centas 3 und der letzten auf Seite 476, welche in der Handschrift mit einer von der hier abgedruckten äusserlich verschiedenen Aufzeichnung der Centas 1 und 2 vereinigt vorkommen, sind die zwölf Abtheilungen statt mit I. Ordo unicus. 1; I. O. 2; I. O. 3; I. O. 4; II. Ordines duo. 1. 1; II. O. 1. 2; II. O. 2. 2; II. O. 3. 3; III. Ordines quatuor. 1. 1. 1. 1; III. O. 1. 1. 2. 2; III. O. 2. 2. 2. 2; IV. Ordines octo 1. 1. 1. 1. 1. 1. 1. 1; hier auf die sonst angewandte Weise mit G. I. 1; G. I. 3; G. I. 5; G. I. 7; G. II. 1; G. II. 2; G. II. 3; G. II. 5; G. IV. 1; G. IV. 2; G. IV. 3; G. VIII. 1; bezeichnet. Die Rechnung ergibt nemlich z. B. 269. I. $3\ [3::4]$; 235. IV. $3\ [6\pdot2::3\pdot5]$; 401. I. $5\ [5::9]$; 577. I. $7\ [7::3]$; 727. II. $5\ [10::3]$. (Genera I statt Genus I auf Seite 469 ist ein Druckfehler).
\par\needspace{5\baselineskip}In Folge von Druckfehlern ist auszulassen:
\begin{longtable}{rcrl}
Centas 27. & G. IV $\dotsfour$ & $16\mathinner{.\,.\,.}$ & $-2624{}^{*}3^{*}$ \\
93 & IV & $16$ & $-9216$ \\
118 & VIII & $4$ & $-11713{}^{*}3^{*}$ \\
\end{longtable}
\par\needspace{5\baselineskip}und hinzuzufügen:
\begin{longtable}{rcrl}
Centas 27. & G. IV $\dotsfour$ & $16\mathinner{.\,.\,.}$ & $-2624{}^{*}2^{*}[16\pdot4::3\pdot16]$ \\
93 & IV & $16$ & $-9216{}^{*}2^{*}[16\pdot4::5\pdot9]$ \\
118 & VIII & $4$ & $-11713{}^{*}2^{*}[4\pdot4\pdot2::31\pdot37\pdot2]$ \\
\end{longtable}
\par\needspace{5\baselineskip}Nach meiner Berechnung ist noch auszulassen:
\begin{longtable}{rcrl}
Centas 10. & G. II $\dotsfour$ & $9\mathinner{.\,.\,.}$ & $-972[6\pdot3::7\pdot13]$ \\
17 & IV & $4$ & $-1660[10\pdot2::11\pdot5]$ \\
20 & IV & $12$ & $-1982[24::3]$ \\
21 & IV & $6$ & $-2096[30\pdot2::3\pdot4]$ \\
23 & IV & $9$ & $-2221[18::10]$ \\
24 & IV & $12$ & $-2376[12\pdot2\pdot2::5\pdot8\pdot8]$ \\
29 & IV & $9$ & $-2887[25::8]$ \\
61 & IV & $7$ & $-6028[12\pdot2::13\pdot4]$ \\
96 & VIII & $13$ & $-9594[20\pdot2\pdot2::31\pdot2\pdot13]$ \\
118 & IV & $25$ & $-11780[16\pdot4\pdot2::3\pdot8\pdot19]$ \\
118 & VIII & $16$ & $-11780[16\pdot4\pdot2::3\pdot8\pdot19]$ \\
119 & VIII & $16$ & $-11840[24\pdot2\pdot2::5\pdot9\pdot7]$ \\
\noalign{\SourceLeaf{533}}
Millias I & G. II $\mathinner{.\,.\,.}$ & $3\mathinner{.\,.\,.}$ & $-541[10::11]$ \\
I & II & $4$ & $-415[10::13]$ \\
I & II & $8$ & $-527[18::3]$ \\
I & II & $8$ & $-722[18::3]$ \\
I & II & $9$ & $-194[20::5]$ \\
I & II & $9$ & $-459[6\pdot3::5\pdot9]$ \\
I & II & $9$ & $-972[6\pdot3::7\pdot13]$ \\
I & II & $11$ & $-842[26::13]$ \\
I & IV & $3$ & $-784[8\pdot2::5\pdot4]$ \\
I & IV & $4$ & $-532[4\pdot2::13\pdot7]$ \\
I & IV & $5$ & $-425[12\pdot2::3\pdot17]$ \\
I & IV & $5$ & $-608[12\pdot2::13\pdot27]$ \\
I & IV & $5$ & $-629[18\pdot2::5\pdot2]$ \\
III & II & $15$ & $-2578[16::13]$ \\
X & I & $111$ & $-9059[117::5]$ \\
\multicolumn{2}{l}{formae $-(15n+13)$ IV} & $4$ & $-2788{}^{*}2^{*}[8\pdot2::19\pdot17]$ \\
\end{longtable}
\par\needspace{5\baselineskip}und hinzuzufügen:
\begin{longtable}{rcrl}
Centas 10. & G. II $\dotsfour$ & $9$ & $-972{}^{*}3^{*}[6\pdot3::7\pdot13]$ \\
17 & IV & $12$ & $-1700[24\pdot2::3\pdot17]$ \\
20 & IV & $12$ & $-1937[24\pdot2::7\pdot2]$ \\
21 & IV & $6$ & $-2097[12\pdot2::47\pdot2]$ \\
23 & IV & $9$ & $-2224[18\pdot2::5\pdot16]$ \\
24 & IV & $12$ & $-2366[24\pdot2::3\pdot2]$ \\
29 & IV & $9$ & $-2885[18\pdot2::3\pdot5]$ \\
61 & IV & $6$ & $-6028[12\pdot2::13\pdot4]$ \\
96 & VIII & $13$ & $-9546[26\pdot2\pdot2::5\pdot3\pdot37]$ \\
118 & IV & $25$ & $-11750[50\pdot2::3\pdot47]$ \\
118 & VIII & $16$ & $-11780{}^{*}2^{*}[16\pdot4\pdot2::3\pdot8\pdot19]$ \\
119 & VIII & $16$ & $-11840{}^{*}2^{*}[16\pdot4\pdot2::3\pdot16\pdot5]$ \\
Millias I. & G. II $\dotsfour$ & $5\mathinner{.\,.\,.}$ & $-415[10::13]$ \\
I & II & $5$ & $-541[10::11]$ \\
I & II & $9$ & $-459{}^{*}3^{*}[6\pdot3::5\pdot9]$ \\
I & II & $9$ & $-527[18::3]$ \\
I & II & $9$ & $-722[18::3]$ \\
I & II & $9$ & $-972{}^{*}3^{*}[6\pdot3::7\pdot13]$ \\
I & II & $10$ & $-194[20::5]$ \\
I & II & $13$ & $-842[26::13]$ \\
I & IV & $2$ & $-532[4\pdot2::13\pdot7]$ \\
I & IV & $4$ & $-784[8\pdot2::5\pdot4]$ \\
I & IV & $6$ & $-425[12\pdot2::3\pdot17]$ \\
I & IV & $6$ & $-608[12\pdot2::13\pdot27]$ \\
I & IV & $9$ & $-629[18\pdot2::5\pdot2]$ \\
III & II & $15$ & $-2518[30::19]$ \\
X & I & $117$ & $-9059[117::5]$ \\
\multicolumn{2}{l}{formae $-(15n+13)$ IV} & $4$ & $-2788[8\pdot2::19\pdot17]$ \\
\end{longtable}
\par Die Tafeln \textit{zur Cyklotechnie} geben für 2452 Zahlen von der Form $aa+1$, $aa+4$, $aa+9$, $\dotsfour aa+81$ die sämmtlichen ungeraden Primtheiler $p$ neben dem zugehörigen $a$ und zwar in solchen Fällen, wo die Primtheiler alle unter 200 liegen, nur dann werden $aa+1$ u. s. f. zerlegbar genannt.
\par Zur leichtern Uebersicht beim Gebrauche hat \textsc{Gauss} für jede Tafel, aus der sich die vollständigen Zerlegungen von Zahlen einer der besonderen Formen bestimmen lassen, eine Hülfstafel aufgestellt, die neben jeder Primzahl $p$ solche Zahlen $a$ enthält, deren um 1 oder $4\mathinner{.\,.\,.}$ vermehrtes Quadrat die Zahl $p$ zum grössten Primtheiler hat.
\par Der Hauptzweck der Tafeln ist die Erleichterung, die sie für die genaue Berechnung der Bögen gewähren, deren Cotangenten gegebene rationale Zahlen sind. Zunächst können nemlich mit ihrer Hülfe die Bögen für kleine Cotangenten aus den Bögen für grosse Cotangenten zusammengesetzt und dadurch die noch erforderlichen Berechnungen der Reihen, welche die Bögen in ihren Cotangenten ausdrücken, auf ein sehr geringes Maass beschränkt werden. Die hierauf hinzielenden Entwickelungen, die sich in dem handschriftlichen Nachlass finden, sind wenig ausgedehnt, die folgende ist die am weitesten fortgeführte. Es bezeichnen darin
\[
\begin{array}{l}
[2]\ [5]\ [13]\ [17]\ [29]\ [37]\ [41]\ [53]\ [61]\mathinner{.\,.\,.}[197]\\
(18)\ (57)\ (239)\ \left(\frac{79}{3}\right)\mathinner{.\,.\,.}
\end{array}
\]
die Bögen der Cotangenten
\[
1\quad2\quad\frac{3}{2}\quad4\quad\frac{5}{2}\quad6\quad\frac{5}{4}\quad\frac{7}{2}\quad\frac{6}{5}\mathinner{.\,.\,.}14\quad18\quad57\quad239\quad\frac{79}{3}\mathinner{.\,.\,.}
\]
Mit Hülfe der Tafeln ist durch Zerlegung von $18+i$, $57+i$, $239+i$ in ihre complexe Primfactoren
\[
\begin{aligned}
(18)&=2[2]-2[5]-[13]\\
(57)&=-[2]+3[5]-[13]\\
(239)&=3[2]\phantom{-2[5]}-4[13]
\end{aligned}
\]
